\documentclass[11pt,a4paper]{beamer}
\usetheme{Madrid}
\usepackage[utf8]{inputenc}
\usepackage{amsmath}
\usepackage{amsfonts}
\usepackage{amssymb}
\usepackage{graphicx}
\usepackage{xcolor}
\usepackage{tikz}
\usepackage{booktabs}
\usepackage{array}
\usepackage{colortbl}
\usepackage{enumitem}

\title{CAMS Exam Preparation}
\subtitle{Domain 6: Building an Anti-Financial Crime Compliance Program (30\% of Exam)}
\author{Based on CAMS Candidate Handbook \& Candidate Feedback}
\date{}

\setbeamercovered{transparent}
\setbeamertemplate{navigation symbols}{}

% Custom highlight commands
\newcommand{\redflag}[1]{\textcolor{red}{\textbf{#1}}}
\newcommand{\bluenote}[1]{\textcolor{blue}{\textbf{#1}}}
\newcommand{\greennote}[1]{\textcolor{green!50!black}{\textbf{#1}}}
\newcommand{\examfocus}[1]{\textcolor{orange!80!black}{\textbf{Exam Focus: #1}}}
\newcommand{\trap}[1]{\textcolor{purple}{\textbf{⚠️ Distractor Trap: #1}}}
\newcommand{\correct}[1]{\textcolor{green!50!black}{\textbf{✓ Correct Answer Approach: #1}}}
\newcommand{\scenario}[1]{\textcolor{blue!70!black}{\textbf{Scenario: #1}}}

\begin{document}
	
	% ============================================================
	% TITLE SLIDE
	% ============================================================
	\begin{frame}
		\titlepage
		\centering
		\textit{Domain 6: 30\% of CAMS Exam | Governance, Controls, Investigations \& Testing}
	\end{frame}
	
	% ============================================================
	% SECTION 1: THREE LINES OF DEFENSE (Governance)
	% ============================================================
	
	\begin{frame}{Section 6.1: The Three Lines of Defense Model}
		\begin{block}{Core Concept}
			A governance framework ensuring separation between risk-taking, risk oversight, and independent assurance.
		\end{block}
		\begin{tabular}{p{0.2\textwidth}|p{0.35\textwidth}|p{0.35\textwidth}}
			\hline
			\rowcolor{lightgray}
			\textbf{Line} & \textbf{Who} & \textbf{Primary Role} \\
			\hline
			\textbf{1st Line} & Business ops, RMs, tellers & Own and manage risk; identify anomalies \\
			\hline
			\textbf{2nd Line} & Compliance, BSA Officer & Oversee risk; set rules; monitor Line 1 \\
			\hline
			\textbf{3rd Line} & Internal Audit & Independent assurance; audit Lines 1 \& 2 \\
			\hline
		\end{tabular}
		\vspace{0.3cm}
		\redflag{Golden Rule:} No single line should both DESIGN and TEST the same controls.
	\end{frame}
	
	\begin{frame}{6.1.1: First Line Violation - The RM Who Ignores Red Flags}
		\examfocus{Common Question Type: Relationship Manager Failure to Escalate}
		
		\scenario{Typical Exam Scenario:}
		A relationship manager at a commercial bank notices that a long-standing corporate client has suddenly started making weekly cash deposits of \$9,500—just below the CTR threshold. The client owns a small retail business. The RM tells the compliance analyst: "I've known this customer for 10 years. He's a good guy. Just close the alert." The analyst closes the alert without further review.
		
		\textbf{What the Exam Tests:} Whether you understand that the RM is First Line (identify anomalies) but the compliance analyst is Second Line (must investigate independently).
		
		\textbf{How to Analyze:}
		\begin{enumerate}
			\item Identify who owns what responsibility: RM flags → compliance investigates
			\item The RM's personal knowledge is NOT a substitute for documented investigation
			\item Closing an alert based on a verbal assurance violates AML program requirements
		\end{enumerate}
		
		\trap{Common Distractor:} "The RM's long relationship with the customer is sufficient evidence to close the alert."
		\correct{Correct Answer Approach:} The compliance analyst must conduct an independent investigation and document findings. Personal assurances from the First Line are not sufficient.
	\end{frame}
	
	\begin{frame}{6.1.2: Second Line Failure - The Analyst Who Closes Alerts Without Documentation}
		\examfocus{Common Question Type: Inadequate Alert Adjudication}
		
		\scenario{Typical Exam Scenario:}
		A transaction monitoring analyst reviews an alert for a customer who made 15 cash deposits of \$8,500-\$9,800 over 30 days. The customer is a restaurant owner. The analyst writes in the case file: "Restaurant business. Cash is normal. No SAR filed." The analyst closes the alert without requesting supporting documentation (daily sales logs, deposit receipts) and without consulting the AML supervisor.
		
		\textbf{What the Exam Tests:} The Second Line's duty to conduct meaningful investigations with documented support.
		
		\textbf{How to Analyze:}
		\begin{enumerate}
			\item Recognize the pattern: multiple sub-\$10,000 deposits = potential structuring
			\item The analyst's assumption ("cash is normal") is not verified
			\item No supporting documentation = audit trail failure
			\item No supervisor consultation = governance failure
		\end{enumerate}
		
		\trap{Common Distractor:} "The analyst correctly applied professional judgment to close the alert."
		\correct{Correct Answer Approach:} The analyst must obtain and review source documents, compare deposits against expected business activity, and document the investigation. Closing an alert based on unverified assumptions is a material deficiency.
	\end{frame}
	
	\begin{frame}{6.1.3: Third Line Independence Violation - Audit Designing Controls}
		\examfocus{Common Question Type: Internal Audit Overstepping Its Role}
		
		\scenario{Typical Exam Scenario:}
		A bank's compliance department lacks the technical expertise to build new transaction monitoring scenarios. The CCO asks the Internal Audit department for help. Audit agrees and spends three months co-designing 14 new scenarios with the compliance team. Six months later, the same audit team conducts the annual independent test of the AML program and issues a report concluding all 14 new scenarios are effective.
		
		\textbf{What the Exam Tests:} The Third Line's independence requirement.
		
		\textbf{How to Analyze:}
		\begin{enumerate}
			\item Internal Audit's role is to TEST controls, not DESIGN them
			\item Designing controls makes Audit part of the Second Line
			\item An auditor cannot objectively test their own work
			\item The audit findings are unreliable due to lack of independence
		\end{enumerate}
		
		\trap{Common Distractor:} "The scenarios were effective, so there is no violation." OR "The CCO appropriately sought help from available resources."
		\correct{Correct Answer Approach:} Internal Audit must remain independent. Designing controls compromises that independence. The bank should hire external consultants or use a different internal team that had no design involvement.
	\end{frame}
	
	\begin{frame}{6.1.4: Board Oversight Violation - Senior Management Hiding Deficiencies}
		\examfocus{Common Question Type: Board's Right to Know}
		
		\scenario{Typical Exam Scenario:}
		A compliance officer discovers that the bank's transaction monitoring system has failed to process 34\% of international wire transfers for 14 months due to a technical misconfiguration. Approximately \$340 million in transactions were never monitored. The compliance officer reports this to senior management. Senior management decides not to inform the Board because "the issue is operational and will be fixed within 90 days."
		
		\textbf{What the Exam Tests:} The Board's oversight role and management's obligation to report material deficiencies.
		
		\textbf{How to Analyze:}
		\begin{enumerate}
			\item A 14-month monitoring gap affecting \$340 million is MATERIAL
			\item Material deficiencies MUST be reported to the Board
			\item Senior management cannot filter or delay reporting
			\item The Board needs this information to exercise oversight
		\end{enumerate}
		
		\trap{Common Distractor:} "Senior management has the authority to decide what information reaches the Board."
		\correct{Correct Answer Approach:} The Board must be informed of material deficiencies that could impact the AML program's effectiveness. Senior management cannot withhold this information.
	\end{frame}
	
	\begin{frame}{6.1.5: Board vs. Management - Who Approves the AML Program?}
		\examfocus{Common Question Type: Distinguishing Board and Management Responsibilities}
		
		\scenario{Typical Exam Scenario:}
		A bank's CCO drafts a revised AML policy. The policy is reviewed by the Legal department and approved by the CCO. The CCO then implements the policy without presenting it to the Board. The CCO argues that since the policy was legally reviewed and complies with regulations, Board approval is unnecessary.
		
		\textbf{What the Exam Tests:} Whether you know the Board must APPROVE the AML program.
		
		\textbf{How to Analyze:}
		\begin{enumerate}
			\item Under BSA/AML requirements, the Board must approve the AML program
			\item Material changes to the program also require Board approval
			\item Management implements; Board oversees and approves
		\end{enumerate}
		
		\trap{Common Distractor:} "The CCO is the designated AML officer with authority to implement policies."
		\correct{Correct Answer Approach:} The Board must approve the AML compliance program and any material changes. The CCO implements, but does not have final approval authority.
	\end{frame}
	
	% ============================================================
	% SECTION 2: INDEPENDENT TESTING (Audit)
	% ============================================================
	
	\begin{frame}{Section 6.2: Independent Testing - What Must Be Tested}
		\begin{block}{Required Scope of Independent Testing}
			The audit must test the \textbf{operational effectiveness} of:
		\end{block}
		\begin{itemize}
			\item Transaction monitoring system (alerts, thresholds, false negative rate)
			\item Sanctions screening implementation (fuzzy logic, name matching)
			\item CDD/KYC documentation quality and completeness
			\item SAR filing timeliness and quality
			\item Training program effectiveness
			\item Governance and oversight structures
		\end{itemize}
		\redflag{Not sufficient:} Only reviewing policies without testing transactions. The audit must test actual samples.
	\end{frame}
	
	\begin{frame}{6.2.1: Testing Transaction Monitoring - The False Negative Problem}
		\examfocus{Common Question Type: Audit Discovers High False Negative Rate}
		
		\scenario{Typical Exam Scenario:}
		An independent auditor tests the bank's transaction monitoring system by running 500 historical SAR cases through the current system. The results show that 387 of the 500 SAR cases (77\%) would NOT have generated an alert under the current thresholds. The compliance director argues this is acceptable because the bank already reviews 22,000 alerts per month and cannot handle more.
		
		\textbf{What the Exam Tests:} Whether you know that false negatives (missed suspicious activity) are a material deficiency.
		
		\textbf{How to Analyze:}
		\begin{enumerate}
			\item A 77\% false negative rate on known SAR cases is EXTREMELY high
			\item This means the system is missing most suspicious activity
			\item Operational burden is NOT a valid excuse for ineffective monitoring
			\item The system must be recalibrated to reduce false negatives
		\end{enumerate}
		
		\trap{Common Distractor:} "High alert volume justifies maintaining current thresholds to control operational costs."
		\correct{Correct Answer Approach:} The monitoring system must be effective first. The bank must document the deficiency, create a remediation plan, and recalibrate thresholds to reduce false negatives.
	\end{frame}
	
	\begin{frame}{6.2.2: Who Can Perform Independent Testing?}
		\begin{block}{Qualification Requirements}
			The tester must have:
		\end{block}
		\begin{enumerate}
			\item \textbf{Independence} - No involvement in designing or operating the controls being tested
			\item \textbf{Expertise} - Sufficient AML/CFT knowledge and qualifications
			\item \textbf{Objectivity} - No conflicts of interest
			\item \textbf{Reporting line} - Reports to Board or Audit Committee (not just CCO)
		\end{enumerate}
		\begin{block}{Acceptable Testers}
			\begin{itemize}
				\item Internal audit (if independent of tested functions)
				\item External auditors
				\item Qualified third-party consultants
			\end{itemize}
		\end{block}
		\redflag{Not Acceptable:} The compliance department testing itself, or audit testing controls they helped design.
	\end{frame}
	
	\begin{frame}{6.2.3: The Prior Involvement Trap}
		\examfocus{Common Question Type: Tester Who Previously Worked in Compliance}
		
		\scenario{Typical Exam Scenario:}
		A bank hires an external consulting firm to perform its independent AML audit. The consulting firm assigns a team that includes an individual who was the bank's BSA Officer until 8 months ago. The individual personally designed several of the policies and procedures now being audited. The bank's CCO sees no issue because the consultant is external.
		
		\textbf{What the Exam Tests:} Whether independence is about the PERSON, not just the firm.
		
		\textbf{How to Analyze:}
		\begin{enumerate}
			\item Independence is individual, not just organizational
			\item A person who designed controls cannot audit those controls
			\item The 8-month separation does not restore objectivity
			\item The consulting firm should assign different personnel
		\end{enumerate}
		
		\trap{Common Distractor:} "The consultant is external, so independence is automatically satisfied."
		\correct{Correct Answer Approach:} The individual tester must have no prior involvement in designing or operating the controls being tested. The consulting firm should assign a different team member.
	\end{frame}
	
	\begin{frame}{6.2.4: Audit Reporting Line - Why It Matters}
		\examfocus{Common Question Type: Audit Reports to CCO}
		
		\scenario{Typical Exam Scenario:}
		A bank's internal audit department completes its annual AML audit and issues a draft report identifying several deficiencies. The audit manager provides the draft to the CCO for review and approval before finalizing. The CCO requests that several critical findings be removed. The audit manager agrees to remove them. The final report shows no deficiencies.
		
		\textbf{What the Exam Tests:} The importance of audit independence and direct Board reporting.
		
		\textbf{How to Analyze:}
		\begin{enumerate}
			\item Audit must report findings directly to the Board/Audit Committee
			\item Management (including CCO) cannot approve or edit audit findings
			\item Allowing management to remove findings destroys audit independence
			\item The final report is unreliable
		\end{enumerate}
		
		\trap{Common Distractor:} "The CCO is responsible for the AML program and should review audit findings for accuracy."
		\correct{Correct Answer Approach:} Audit findings must be reported directly to the Board/Audit Committee without management filtering or approval. The CCO may receive a copy but cannot control the content.
	\end{frame}
	
	% ============================================================
	% SECTION 3: TRANSACTION MONITORING CALIBRATION
	% ============================================================
	
	\begin{frame}{Section 6.3: Transaction Monitoring - False Positives vs. False Negatives}
		\begin{block}{The Core Trade-Off}
			Every AML monitoring system faces this calibration decision.
		\end{block}
		\begin{tabular}{p{0.3\textwidth}|p{0.6\textwidth}}
			\hline
			\rowcolor{lightgray}
			\textbf{Term} & \textbf{Definition and Risk} \\
			\hline
			\textbf{False Positive} & Legitimate activity triggers an alert. \redflag{Risk:} Wasted analyst time, operational inefficiency. \\
			\hline
			\textbf{False Negative} & Suspicious activity fails to generate an alert. \redflag{Risk:} MISSED criminal activity. Direct AML program deficiency. \\
			\hline
		\end{tabular}
		\vspace{0.3cm}
		\redflag{Which is worse? FALSE NEGATIVES.} Missing actual crime is always worse than reviewing too many alerts.
		\bluenote{Regulator Priority:} Detection effectiveness > operational efficiency.
	\end{frame}
	
	\begin{frame}{6.3.1: The Over-Filtering Trap (Too Many False Negatives)}
		\examfocus{Common Question Type: Director Raises Thresholds to Reduce Alerts}
		
		\scenario{Typical Exam Scenario:}
		A bank generates 8,000 alerts per month. The compliance director reviews metrics showing that 96\% of alerts are closed as false positives. Under pressure to reduce operational costs, the director orders the analytics team to raise all threshold values by 250\%. The team leader warns this will increase false negatives. The director replies: "We can't afford to review 8,000 alerts. Any missed activity is acceptable as long as we still file some SARs."
		
		\textbf{What the Exam Tests:} Whether you know that over-filtering (raising thresholds too high) creates unacceptable false negative risk.
		
		\textbf{How to Analyze:}
		\begin{enumerate}
			\item 250\% threshold increase is extreme and arbitrary
			\item Many suspicious transactions will fall below new thresholds
			\item The director is prioritizing cost over effectiveness
			\item Regulators will view this as a material deficiency
		\end{enumerate}
		
		\trap{Common Distractor:} "The director is correctly managing operational efficiency under the risk-based approach."
		\correct{Correct Answer Approach:} Raising thresholds must be based on data analysis, not arbitrary cost-cutting. The bank must validate that new thresholds still detect known typologies. Effectiveness cannot be sacrificed for efficiency.
	\end{frame}
	
	\begin{frame}{6.3.2: The Under-Filtering Trap (Too Many False Positives)}
		\examfocus{Common Question Type: System Generates Too Many Alerts to Investigate}
		
		\scenario{Typical Exam Scenario:}
		A bank sets its thresholds very low to ensure no suspicious activity is missed. The system generates 50,000 alerts per month. The compliance team has only 5 analysts, each capable of reviewing 200 alerts per month. The backlog grows to 40,000 uninvestigated alerts. Most alerts are never reviewed because the team cannot keep up.
		
		\textbf{What the Exam Tests:} Whether you know that uninvestigated alerts are as bad as no alerts at all.
		
		\textbf{How to Analyze:}
		\begin{enumerate}
			\item Alerts that are never reviewed provide no value
			\item The system is generating more alerts than the team can handle
			\item The bank must either increase staffing OR recalibrate thresholds
			\item Both under-filtering AND over-filtering are problems
		\end{enumerate}
		
		\trap{Common Distractor:} "Lower thresholds are always better because they catch more activity."
		\correct{Correct Answer Approach:} The system must be calibrated so that the team can reasonably review all alerts. Unreviewed alerts are a control failure. The bank needs the right balance of thresholds AND adequate staffing.
	\end{frame}
	
	\begin{frame}{6.3.3: Model Drift - Why Monitoring Systems Degrade Over Time}
		\examfocus{Common Question Type: Model Performance Degradation}
		
		\scenario{Typical Exam Scenario:}
		A bank deployed a machine learning monitoring model in January 2023. The validation report at deployment showed 72\% precision and 81\% recall. In June 2025, a validation review finds precision has dropped to 51\% and recall to 62\%. The bank introduced new payment products, customer transaction volumes increased 180\%, and criminal typologies have shifted from cash structuring to crypto layering. The model was never recalibrated after deployment.
		
		\textbf{What the Exam Tests:} Model drift and the need for regular recalibration.
		
		\textbf{How to Analyze:}
		\begin{enumerate}
			\item Models degrade when customer behavior, products, or typologies change
			\item The bank failed to establish a refresh cadence
			\item Annual recalibration or trigger-based retraining is required
			\item The performance degradation requires immediate remediation
		\end{enumerate}
		
		\trap{Common Distractor:} "The model still has recall above 60\%, which is acceptable."
		\correct{Correct Answer Approach:} The bank must establish model monitoring and refresh processes. Significant performance degradation requires recalibration or model replacement.
	\end{frame}
	
	\begin{frame}{6.3.4: Scenario Validation Frequency - Annual vs. Trigger-Based}
		\examfocus{Common Question Type: Missing Interim Validation After Material Changes}
		
		\scenario{Typical Exam Scenario:}
		A bank's AML policy requires full scenario validation "at least annually." Eight months after the last validation, the bank introduces cryptocurrency custody services, acquires a competitor adding 35\% more customers, and sees a 400\% increase in SARs for human trafficking typologies. The bank does not conduct any interim validation. A regulator asks why.
		
		\textbf{What the Exam Tests:} The need for event-driven validation outside the annual cycle.
		
		\textbf{How to Analyze:}
		\begin{enumerate}
			\item Material changes trigger the need for interim validation
			\item New products, customer growth, and emerging typologies are all trigger events
			\item Annual validation is a minimum, not a maximum
			\item The bank should have trigger-based validation in its policy
		\end{enumerate}
		
		\trap{Common Distractor:} "Annual validation is sufficient because the policy only requires annual testing."
		\correct{Correct Answer Approach:} Material changes to risk profile, products, customer composition, or typologies should trigger interim validation. The policy should define specific trigger events.
	\end{frame}
	
	% ============================================================
	% SECTION 4: DATA QUALITY AND RECONCILIATION
	% ============================================================
	
	\begin{frame}{Section 6.4: Data Quality - The Silent Failure}
		\begin{block}{Core Concept}
			A monitoring system is only as good as the data it receives. Missing data = missing alerts.
		\end{block}
		\textbf{Critical Controls:}
		\begin{itemize}
			\item Reconciliation between source systems and the monitoring system
			\item Data completeness and accuracy checks
			\item Error handling and alerting for failed data feeds
			\item Regular validation that all required transactions are ingested
		\end{itemize}
		\redflag{The worst data failure:} Silent failure where data is missing but no error messages alert the team.
	\end{frame}
	
	\begin{frame}{6.4.1: The Silent Data Feed Failure}
		\examfocus{Common Question Type: Undetected Data Gap}
		
		\scenario{Typical Exam Scenario:}
		A bank's transaction monitoring system relies on daily data feeds from 17 upstream systems. An audit reveals that for 14 months, the feed from the international wire system silently failed for 34\% of all wires. Those transactions never entered the monitoring system, never generated alerts, and were never reviewed. An estimated \$340 million in wires were not monitored. The failure was a technical misconfiguration with no error messages. No reconciliation controls existed to detect the gap.
		
		\textbf{What the Exam Tests:} The importance of data reconciliation controls.
		
		\textbf{How to Analyze:}
		\begin{enumerate}
			\item 34\% missing data = massive monitoring gap
			\item No alerts could be generated for transactions the system never received
			\item The absence of reconciliation controls is a foundational failure
			\item The bank must remediate AND review missed transactions retrospectively
		\end{enumerate}
		
		\trap{Common Distractor:} "The misconfiguration was accidental, so no compliance violation occurred."
		\correct{Correct Answer Approach:} The bank must have controls to verify complete data ingestion. The 14-month gap is a material deficiency requiring retrospective review and SAR filing for identified suspicious activity.
	\end{frame}
	
	\begin{frame}{6.4.2: Missing Customer Identification - The Unlinked Transaction Problem}
		\examfocus{Common Question Type: Transactions That Cannot Be Linked to Customers}
		
		\scenario{Typical Exam Scenario:}
		A bank's transaction monitoring system flags a series of large wire transfers. When the analyst tries to investigate, the system shows the transactions are associated with a "Customer ID: NULL" value. The data issue has been present for 8 months. The bank's IT team explains that when customers open accounts through the mobile app, the unique identifier sometimes fails to populate. Approximately 12\% of transactions cannot be linked to any customer record.
		
		\textbf{What the Exam Tests:} The importance of customer identification in transaction monitoring.
		
		\textbf{How to Analyze:}
		\begin{enumerate}
			\item Unlinked transactions cannot be investigated properly
			\item The bank cannot determine customer risk or conduct CDD
			\item This is a data quality failure with direct AML impact
			\item The bank cannot file accurate SARs without customer identification
		\end{enumerate}
		
		\trap{Common Distractor:} "IT issues are operational and not a compliance concern."
		\correct{Correct Answer Approach:} Missing customer identification data is a material AML deficiency. The bank must fix the root cause and develop a process to identify and investigate orphaned transactions.
	\end{frame}
	
	% ============================================================
	% SECTION 5: ALERT ADJUDICATION AND INVESTIGATION
	% ============================================================
	
	\begin{frame}{Section 6.5: Alert Adjudication - The Investigator Variation Problem}
		\begin{block}{Core Concept}
			Different analysts should reach similar conclusions on similar facts. Wide variation indicates lack of standards.
		\end{block}
		\textbf{Required Controls:}
		\begin{itemize}
			\item Written adjudication guidelines for each scenario
			\item Required documentation standards
			\item Quality assurance reviews of closed alerts
			\item Calibration sessions to align analyst judgment
		\end{itemize}
		\redflag{Red Flag:} One analyst closes 95\% of alerts as false positives; another escalates 70\% for the same scenario.
	\end{frame}
	
	\begin{frame}{6.5.1: The Inconsistent Analyst Problem}
		\examfocus{Common Question Type: Wide Variation in Analyst Decisions}
		
		\scenario{Typical Exam Scenario:}
		A bank's QA team reviews 200 closed alerts for a cash deposit scenario. Analyst A closed 94\% as false positives with a note: "Customer is a retailer; cash is normal." Analyst B, reviewing similar accounts for the same scenario, escalated 67\% for further investigation and filed SARs on 22\%. Neither analyst's file includes supporting documentation. The bank has no written adjudication guidelines for this scenario.
		
		\textbf{What the Exam Tests:} The need for standardization and documentation.
		
		\textbf{How to Analyze:}
		\begin{enumerate}
			\item 94\% vs. 67\% variation is extreme and unacceptable
			\item Without written guidelines, analysts apply inconsistent standards
			\item No documentation means the bank cannot demonstrate investigations occurred
			\item The bank needs written guidelines AND QA oversight
		\end{enumerate}
		
		\trap{Common Distractor:} "Different analysts apply professional judgment differently; variation is acceptable."
		\correct{Correct Answer Approach:} The bank must have written adjudication standards, required documentation, and QA processes to ensure consistent, defensible decisions.
	\end{frame}
	
	\begin{frame}{6.5.2: The Uninvestigated Alert - Assuming Without Evidence}
		\examfocus{Common Question Type: Closing Alerts Based on Unverified Assumptions}
		
		\scenario{Typical Exam Scenario:}
		An alert flags a customer who made 23 cash deposits between \$8,400 and \$9,700 over 45 days. The customer is a restaurant owner. The analyst writes: "Restaurant business. Cash from customers. No SAR needed." The analyst closes the alert without requesting daily sales logs, comparing deposits to prior periods, or checking if the deposit pattern changed after a new manager was hired.
		
		\textbf{What the Exam Tests:} The difference between assumption and investigation.
		
		\textbf{How to Analyze:}
		\begin{enumerate}
			\item The deposit pattern is a classic structuring indicator
			\item The analyst assumed legitimacy without verification
			\item No supporting documentation was obtained
			\item The alert was closed without a meaningful investigation
		\end{enumerate}
		
		\trap{Common Distractor:} "The analyst used business context to make a reasonable judgment."
		\correct{Correct Answer Approach:} Analysts must obtain and review supporting documentation. Assumptions without verification do not constitute an investigation. The bank needs documented evidence that a reasonable investigation occurred.
	\end{frame}
	
	% ============================================================
	% SECTION 6: SAR FILING, TIPPING-OFF, AND CONFIDENTIALITY
	% ============================================================
	
	\begin{frame}{Section 6.6: SAR Filing - Timing and Standard}
		\begin{block}{Core Concept}
			SARs must be filed immediately upon suspicion, regardless of transaction value.
		\end{block}
		\textbf{Key Principles:}
		\begin{itemize}
			\item Standard = reasonable suspicion (not certainty, not proof)
			\item No minimum monetary threshold for ML suspicion
			\item Do NOT delay to gather more evidence
			\item Do NOT investigate to the point of certainty before filing
			\item File and continue investigating if needed
		\end{itemize}
		\redflag{Common violation:} "We need more evidence before filing" → This is incorrect. Reasonable suspicion is enough.
	\end{frame}
	
	\begin{frame}{6.6.1: The "Need More Evidence" Trap}
		\examfocus{Common Question Type: Delaying SAR Filing for More Proof}
		
		\scenario{Typical Exam Scenario:}
		A compliance analyst identifies a pattern of transactions that raises reasonable suspicion of money laundering. The analyst's supervisor says: "We don't have enough evidence yet. Keep investigating and file when you have conclusive proof." The analyst continues investigating for 60 days without filing a SAR.
		
		\textbf{What the Exam Tests:} The SAR filing standard is reasonable suspicion, not proof.
		
		\textbf{How to Analyze:}
		\begin{enumerate}
			\item Reasonable suspicion is the legal standard
			\item Conclusive proof is NOT required
			\item Delaying filing is a violation
			\item The bank can (and should) file and continue investigating
		\end{enumerate}
		
		\trap{Common Distractor:} "The supervisor correctly wants more evidence to ensure the SAR is accurate."
		\correct{Correct Answer Approach:} The SAR must be filed immediately upon reasonable suspicion. The bank can file and then continue investigating. Delaying filing is a regulatory violation.
	\end{frame}
	
	\begin{frame}{6.6.2: Tipping-Off - What Constitutes a Violation}
		\examfocus{Common Question Type: Identifying Tipping-Off Scenarios}
		
		\scenario{Typical Exam Scenario (Violation):}
		A customer asks a branch teller why a wire transfer is delayed. The teller says: "I'm sorry, but our compliance team is reviewing your account for some unusual activity." The customer immediately becomes agitated and closes the account. No SAR has been filed yet; the investigation is ongoing.
		
		\textbf{What the Exam Tests:} Tipping-off includes disclosing that an investigation EXISTS, not just that a SAR was filed.
		
		\textbf{How to Analyze:}
		\begin{enumerate}
			\item Tipping-off occurs when you disclose an investigation or SAR filing
			\item You do NOT need to mention "SAR" by name to tip off
			\item Telling a customer they are under compliance review IS tipping-off
			\item The teller should have said nothing or given a neutral response
		\end{enumerate}
		
		\trap{Common Distractor:} "No SAR was filed yet, so no tipping-off occurred."
		\correct{Correct Answer Approach:} Tipping-off includes disclosing the existence of an internal investigation. The bank must train staff not to disclose any compliance scrutiny to customers.
	\end{frame}
	
	\begin{frame}{6.6.3: Permissible Disclosures - When Can You Share SAR Information?}
		\examfocus{Common Question Type: Exceptions to Tipping-Off Rules}
		
		\scenario{Typical Exam Scenario (Permissible):}
		Two banks have both filed SARs on different individuals who appear to be part of the same money laundering network. Both banks have filed FinCEN 314(b) opt-in notifications. The BSA Officers of the two banks meet to share information about the suspicious activity.
		
		\textbf{What the Exam Tests:} The 314(b) exception to tipping-off rules.
		
		\textbf{How to Analyze:}
		\begin{enumerate}
			\item 314(b) permits voluntary information sharing between opted-in institutions
			\item Sharing is protected by safe harbor
			\item This is NOT tipping-off because it is explicitly permitted by law
			\item Both banks must have filed the opt-in notification
		\end{enumerate}
		
		\trap{Common Distractor:} "Sharing SAR information with another bank always constitutes tipping-off."
		\correct{Correct Answer Approach:} 314(b) sharing is a lawful exception to tipping-off rules. Banks that have opted in can share information about suspected ML/TF activity.
	\end{frame}
	
	\begin{frame}{6.6.4: Internal Sharing - Legal Department and Need-to-Know}
		\examfocus{Common Question Type: When Internal Sharing Is Permissible}
		
		\scenario{Typical Exam Scenario (Permissible):}
		A compliance analyst has drafted a SAR. Before filing, the analyst shares the draft with the bank's legal department for review. The legal department advises on potential legal risks. The analyst then files the SAR.
		
		\textbf{What the Exam Tests:} Internal sharing on a need-to-know basis is permissible.
		
		\textbf{How to Analyze:}
		\begin{enumerate}
			\item The legal department has a legitimate need to know
			\item Sharing is internal, not external
			\item The SAR remains confidential and is not disclosed to the customer
			\item This is standard and permissible practice
		\end{enumerate}
		
		\trap{Common Distractor:} "Any sharing of SAR information before filing is tipping-off."
		\correct{Correct Answer Approach:} Internal sharing with legal, compliance management, and audit on a need-to-know basis is permissible and does not constitute tipping-off.
	\end{frame}
	
	% ============================================================
	% SECTION 7: RESPONDING TO LAW ENFORCEMENT
	% ============================================================
	
	\begin{frame}{Section 6.7: Responding to Subpoenas and Law Enforcement Requests}
		\begin{block}{Core Concept}
			Financial institutions have a legal obligation to respond to valid law enforcement requests.
		\end{block}
		\textbf{Key Principles:}
		\begin{itemize}
			\item Comply promptly within specified timeframe
			\item Maintain internal record of everything produced
			\item Do NOT notify the customer (unless subpoena requires disclosure or is quashed)
			\item Seek legal advice if uncertain about scope or validity
		\end{itemize}
		\redflag{What NOT to do:} Destroy records, notify the customer to warn them, or refuse to comply without legal basis.
	\end{frame}
	
	\begin{frame}{6.7.1: The Subpoena - Proper Response}
		\examfocus{Common Question Type: Bank Receives Law Enforcement Subpoena}
		
		\scenario{Typical Exam Scenario:}
		A bank receives a federal grand jury subpoena demanding all account records, wire transfer logs, and signature cards for a specific corporate customer under investigation for money laundering. The subpoena includes a non-disclosure provision prohibiting the bank from notifying the customer.
		
		\textbf{What the Exam Tests:} Proper response to a subpoena with a non-disclosure order.
		
		\textbf{How to Analyze:}
		\begin{enumerate}
			\item The bank must comply with the subpoena
			\item The non-disclosure provision means the bank CANNOT tell the customer
			\item The bank must maintain an internal record of everything produced
			\item The bank should not destroy any responsive records
		\end{enumerate}
		
		\trap{Common Distractor:} "The bank should notify the customer so they can respond to the subpoena."
		\correct{Correct Answer Approach:} Comply promptly, maintain internal records, and do not violate the non-disclosure provision by notifying the customer.
	\end{frame}
	
	\begin{frame}{6.7.2: The "Keep Open" Request - Balancing SAR Closure and LE Needs}
		\examfocus{Common Question Type: Law Enforcement Asks Bank to Keep Account Open}
		
		\scenario{Typical Exam Scenario:}
		A bank has filed three SARs on an account over six months and decided to close the relationship. Before closing, a federal agent contacts the bank in writing, requesting that the bank keep the account open for an additional 90 days to allow an ongoing undercover investigation to continue. The agent says closing the account would alert the subjects.
		
		\textbf{What the Exam Tests:} How to handle a law enforcement "keep open" request.
		
		\textbf{How to Analyze:}
		\begin{enumerate}
			\item The bank should generally comply with written LE requests
			\item The request must be documented in writing
			\item Senior management approval should be obtained
			\item The bank should continue enhanced monitoring during the extension
		\end{enumerate}
		
		\trap{Common Distractor:} "The bank must close the account immediately regardless of the LE request to minimize its own risk."
		\correct{Correct Answer Approach:} The bank should document the request in writing, obtain senior management approval, and keep the account open as requested, with continued enhanced monitoring.
	\end{frame}
	
	% ============================================================
	% SECTION 8: DE-RISKING AND FINANCIAL INCLUSION
	% ============================================================
	
	\begin{frame}{Section 6.8: De-Risking - What It Is and Why It's Problematic}
		\begin{block}{Core Concept}
			De-risking = indiscriminate termination of entire categories of customer relationships without individual risk assessment.
		\end{block}
		\textbf{Why It's Problematic:}
		\begin{itemize}
			\item Drives legitimate activity into unregulated, informal channels
			\item Harms financial inclusion (remittance corridors, developing economies)
			\item Violates the risk-based approach (requires individual assessment)
			\item May contradict regulatory expectations
		\end{itemize}
		\redflag{Key distinction:} Individual, documented termination based on risk = permissible. Blanket termination by geography or category = de-risking (problematic).
	\end{frame}
	
	\begin{frame}{6.8.1: Blanket Termination vs. Individual Assessment}
		\examfocus{Common Question Type: Bank Terminates All Relationships in a Region}
		
		\scenario{Typical Exam Scenario:}
		A global bank terminates all correspondent relationships with banks headquartered in Sub-Saharan Africa. The termination applies to all 22 respondent banks in the region, including three that have passed recent FATF mutual evaluations with "Compliant" ratings. The bank conducted no individual risk assessments of the 22 banks. The compliance director states: "The region is high-risk, so we are exiting entirely."
		
		\textbf{What the Exam Tests:} The difference between permissible risk-based exit and impermissible de-risking.
		
		\textbf{How to Analyze:}
		\begin{enumerate}
			\item Blanket termination without individual assessment violates the RBA
			\item Compliant banks are being terminated along with non-compliant ones
			\item This drives legitimate activity out of the formal system
			\item FATF and Wolfsberg guidance discourages this approach
		\end{enumerate}
		
		\trap{Common Distractor:} "Banks have absolute discretion to choose their correspondent partners."
		\correct{Correct Answer Approach:} The risk-based approach requires individual, documented risk assessments. Blanket categorical termination without assessment is de-risking and is discouraged by regulators and FATF.
	\end{frame}
	
	\begin{frame}{6.8.2: Documented Risk-Based Exit - The Permissible Alternative}
		\examfocus{Common Question Type: Properly Documented Termination}
		
		\scenario{Typical Exam Scenario:}
		A bank terminates its correspondent relationship with a specific respondent bank after: (1) conducting an individual risk assessment; (2) identifying that the respondent bank services 14 unlicensed MSBs; (3) issuing two RFIs that went unanswered; (4) documenting the risk assessment and decision in writing. The respondent bank is one of several in its jurisdiction; others were retained.
		
		\textbf{What the Exam Tests:} The difference between de-risking and justified risk-based exit.
		
		\textbf{How to Analyze:}
		\begin{enumerate}
			\item Individual assessment was conducted (not blanket termination)
			\item Specific risk factors were identified
			\item Attempts to remediate (RFIs) were made
			\item Decision was documented
			\item Other banks in the same jurisdiction were retained (not blanket)
		\end{enumerate}
		
		\trap{Common Distractor:} "Any termination of a correspondent relationship is de-risking."
		\correct{Correct Answer Approach:} Justified risk-based exit based on individual documented assessment is permissible and is NOT de-risking. The key factors are individual assessment, documentation, and attempt to remediate.
	\end{frame}
	
	% ============================================================
	% SECTION 9: REQUESTS FOR INFORMATION (RFIs)
	% ============================================================
	
	\begin{frame}{Section 6.9: Requests for Information (RFIs)}
		\begin{block}{Core Concept}
			RFIs are formal requests sent to customers or respondent banks asking for clarification or documentation of suspicious activity.
		\end{block}
		\textbf{Best Practices:}
		\begin{itemize}
			\item Use when you need more information to resolve an alert
			\item Document the RFI and any response (or lack thereof)
			\item Do NOT disclose that an investigation or SAR is pending (tipping-off)
			\item Failure to respond is itself a red flag
		\end{itemize}
		\redflag{Critical:} RFIs must be phrased neutrally. Do NOT say "we are investigating you for money laundering."
	\end{frame}
	
	\begin{frame}{6.9.1: The Neutral RFI vs. Tipping-Off}
		\examfocus{Common Question Type: Proper RFI Wording}
		
		\scenario{Typical Exam Scenario (Violation):}
		A compliance analyst calls a customer and says: "Our AML system flagged your deposits as potential structuring. We need to know where you're getting all this cash, or we might have to file a report." The customer becomes angry and hangs up.
		
		\textbf{What the Exam Tests:} RFIs cannot disclose an investigation or potential SAR filing.
		
		\textbf{How to Analyze:}
		\begin{enumerate}
			\item Mentioning "AML system" and "potential structuring" is tipping-off
			\item Mentioning "file a report" is explicit tipping-off
			\item The customer now knows they are under investigation
			\item The RFI was phrased improperly
		\end{enumerate}
		
		\trap{Common Distractor:} "The analyst was honest with the customer, which is good customer service."
		\correct{Correct Answer Approach:} RFIs must be neutral. Example: "We noticed an increase in cash deposits. Can you provide documentation explaining the source of these funds?" No mention of AML, structuring, or reporting.
	\end{frame}
	
	% ============================================================
	% SECTION 10: AML TRAINING REQUIREMENTS
	% ============================================================
	
	\begin{frame}{Section 6.10: AML Training - Role-Based Requirements}
		\begin{block}{Core Concept}
			Training must be tailored to the specific roles and risks of each employee group.
		\end{block}
		\textbf{Role-Based Training Examples:}
		\begin{itemize}
			\item \textbf{Tellers:} Cash structuring, CTR triggers, identifying suspicious customer behavior
			\item \textbf{Relationship Managers:} PEP identification, EDD requirements, red flags in customer profiles
			\item \textbf{Board of Directors:} Governance oversight, risk strategy, regulatory expectations
			\item \textbf{IT/Technology:} Data security, system controls, audit trail requirements
			\item \textbf{Compliance Staff:} Advanced investigation techniques, SAR writing, regulatory updates
		\end{itemize}
		\redflag{Deficiency:} The same training module delivered to all employees regardless of role.
	\end{frame}
	
	\begin{frame}{6.10.1: The One-Size-Fits-All Training Failure}
		\examfocus{Common Question Type: Ineffective Training Program}
		
		\scenario{Typical Exam Scenario:}
		A bank delivers the same two-hour online AML training module to all 5,000 employees, including tellers, relationship managers, the Board of Directors, and IT staff. The training covers the technical details of transaction monitoring thresholds. The external auditor rates the training program as "ineffective" because Board members do not need to know threshold calculations, and tellers received no training on cash structuring indicators.
		
		\textbf{What the Exam Tests:} Training must be role-specific and relevant.
		
		\textbf{How to Analyze:}
		\begin{enumerate}
			\item Tellers need training on cash structuring and CTRs, not model calibration
			\item Board members need governance training, not technical threshold details
			\item One-size-fits-all training is ineffective
			\item The bank must develop role-specific modules
		\end{enumerate}
		
		\trap{Common Distractor:} "All employees need the same training to ensure consistent institutional coverage."
		\correct{Correct Answer Approach:} Training must be tailored to the specific roles, responsibilities, and ML/TF risks faced by each employee group. Regulators expect role-based training.
	\end{frame}
	
	\begin{frame}{6.10.2: Board Training - Governance, Not Operations}
		\examfocus{Common Question Type: Appropriate Board Training Content}
		
		\scenario{Typical Exam Scenario:}
		A bank's annual AML training for the Board of Directors covers: (1) how to identify suspicious transactions in transaction monitoring alerts; (2) how to file a Currency Transaction Report; (3) the bank's AML risk appetite and governance oversight responsibilities; (4) regulatory expectations for Board oversight.
		
		\textbf{What the Exam Tests:} Board training should focus on governance, not operations.
		
		\textbf{How to Analyze:}
		\begin{enumerate}
			\item Board members do NOT file CTRs or review alerts (operational)
			\item Board training should cover risk appetite, governance, and oversight
			\item Items 1 and 2 are inappropriate for Board training
			\item Items 3 and 4 are appropriate for Board training
		\end{enumerate}
		
		\trap{Common Distractor:} "Board members need to understand all aspects of the AML program, including operations."
		\correct{Correct Answer Approach:} Board training should focus on governance, oversight, risk strategy, and regulatory expectations—not operational tasks like filing CTRs or reviewing alerts.
	\end{frame}
	
	% ============================================================
	% SECTION 11: WHISTLEBLOWER PROTECTIONS
	% ============================================================
	
	\begin{frame}{Section 6.11: Whistleblower Protections}
		\begin{block}{Core Concept}
			Laws protect employees who report AML violations from retaliation.
		\end{block}
		\textbf{Key Protections (U.S.):}
		\begin{itemize}
			\item Bank Secrecy Act (31 U.S.C. § 5328) prohibits retaliation
			\item Dodd-Frank Act Section 922 (SEC-related disclosures)
			\item FIRREA (financial institution anti-retaliation)
			\item File complaint with Department of Labor (OSHA) within 180 days
		\end{itemize}
		\redflag{What constitutes retaliation:} Termination, demotion, suspension, harassment, performance improvement plans (PIPs) with no prior issues.
	\end{frame}
	
	\begin{frame}{6.11.1: The Retaliatory PIP}
		\examfocus{Common Question Type: Identifying Whistleblower Retaliation}
		
		\scenario{Typical Exam Scenario:}
		A compliance analyst reports internally that senior management is suppressing SAR filings on a high-revenue corporate client. Three weeks later, the analyst receives a performance improvement plan (PIP) citing "communication issues"—her first negative review in eight years. She believes the PIP is retaliatory.
		
		\textbf{What the Exam Tests:} Recognizing retaliation and knowing the protected avenues.
		
		\textbf{How to Analyze:}
		\begin{enumerate}
			\item Timing (3 weeks after report) suggests retaliation
			\item First negative review in 8 years is suspicious
			\item The BSA prohibits retaliation for reporting BSA violations
			\item She can file a complaint with DOL/OSHA within 180 days
		\end{enumerate}
		
		\trap{Common Distractor:} "Internal reports are not protected; only external reports to FinCEN are protected."
		\correct{Correct Answer Approach:} The BSA prohibits retaliation for internal reports as well. The analyst should document everything and file a complaint with DOL/OSHA.
	\end{frame}
	
	% ============================================================
	% SECTION 12: FUNCTIONAL DEPARTMENTS (Data Privacy, Cybersecurity)
	% ============================================================
	
	\begin{frame}{Section 6.12: Supporting Departments - Data Privacy and Cybersecurity}
		\begin{block}{Core Concept}
			AML compliance does not operate in isolation. Other departments enable effective AFC programs.
		\end{block}
		\textbf{Key Supporting Functions:}
		\begin{itemize}
			\item \textbf{Data Privacy:} Ensures customer data is protected while enabling lawful information sharing
			\item \textbf{Cybersecurity:} Protects AML systems from attack; identifies cyber-related red flags (ransomware, unauthorized access)
			\item \textbf{IT/Technology:} Maintains transaction monitoring systems, data feeds, and audit trails
			\item \textbf{Human Resources:} Employee due diligence, insider threat monitoring
			\item \textbf{Vendor Management:} Third-party due diligence for outsourced AML functions
		\end{itemize}
	\end{frame}
	
	\begin{frame}{6.12.1: Data Privacy vs. AML - The GDPR Exception}
		\examfocus{Common Question Type: Resolving Privacy vs. AML Conflicts}
		
		\scenario{Typical Exam Scenario:}
		A European bank receives a GDPR "right to erasure" request from a customer demanding deletion of all personal data, including KYC documents and transaction records. The customer was the subject of a SAR filed 18 months ago. The Data Privacy Officer argues GDPR requires deletion. The AML Officer argues AML recordkeeping requires retention.
		
		\textbf{What the Exam Tests:} GDPR's exception for legal obligations (AML retention).
		
		\textbf{How to Analyze:}
		\begin{enumerate}
			\item GDPR Article 17(3) exempts data retention required by law
			\item AML recordkeeping (minimum 5 years) is a legal obligation
			\item The bank can lawfully refuse to delete AML-required records
			\item SAR confidentiality also prohibits acknowledging the SAR to the customer
		\end{enumerate}
		
		\trap{Common Distractor:} "GDPR's right to erasure is absolute and overrides all other laws."
		\correct{Correct Answer Approach:} AML retention requirements are a legal exception to GDPR erasure rights. The bank may refuse deletion for records it is required to retain.
	\end{frame}
	
	\begin{frame}{6.12.2: Cybersecurity and AML - The Ransomware Red Flag}
		\examfocus{Common Question Type: Cyber-Related Red Flags}
		
		\scenario{Typical Exam Scenario:}
		A bank's cybersecurity team detects that a corporate customer's systems have been compromised by ransomware. The customer pays the ransom in Bitcoin to a wallet address that appears on OFAC's SDN list. The cybersecurity team notifies the AML team.
		
		\textbf{What the Exam Tests:} Interplay between cybersecurity and AML.
		
		\textbf{How to Analyze:}
		\begin{enumerate}
			\item Paying ransom to an SDN-listed wallet is an OFAC violation
			\item The transaction may also constitute money laundering (proceeds of cybercrime)
			\item The bank should file a SAR
			\item The bank should consider whether to block the transaction
		\end{enumerate}
		
		\trap{Common Distractor:} "The bank should process the payment to avoid harm to the customer."
		\correct{Correct Answer Approach:} AML and cybersecurity teams must coordinate. Paying ransom to an SDN-listed wallet requires sanctions review and potential SAR filing.
	\end{frame}
	
	% ============================================================
	% SECTION 13: EMPLOYEE DUE DILIGENCE AND INSIDER THREATS
	% ============================================================
	
	\begin{frame}{Section 6.13: Employee Due Diligence and Insider Threats}
		\begin{block}{Core Concept}
			Insider threats (employees misusing their access) are a significant AML risk.
		\end{block}
		\textbf{Key Controls:}
		\begin{itemize}
			\item Background checks on hiring
			\item Ongoing employee monitoring (not just customer monitoring)
			\item Segregation of duties (no single employee should control end-to-end processes)
			\item Access controls and audit trails
			\item Mandatory vacation policies (to detect ongoing fraud)
		\end{itemize}
		\redflag{Red flags:} Employee living beyond means, refusing vacation, accessing accounts outside job scope.
	\end{frame}
	
	\begin{frame}{6.13.1: The Insider Threat - Employee Override}
		\examfocus{Common Question Type: Identifying Insider Threat Indicators}
		
		\scenario{Typical Exam Scenario:}
		An AML analyst notices that a senior relationship manager has overridden 23 transaction monitoring alerts in the past year for the same corporate client. Each override note states: "Client is legitimate." The relationship manager has access to the override function and does not require supervisory approval for overrides. The relationship manager has recently purchased a new luxury car and vacation home.
		
		\textbf{What the Exam Tests:} Insider threat red flags.
		
		\textbf{How to Analyze:}
		\begin{enumerate}
			\item 23 overrides for the same client is unusual
			\item No supervisory approval for overrides is a control weakness
			\item Lifestyle changes may indicate illicit enrichment
			\item The bank should investigate the relationship manager
		\end{enumerate}
		
		\trap{Common Distractor:} "The relationship manager has authority to override alerts."
		\correct{Correct Answer Approach:} Multiple overrides for the same client, combined with lifestyle changes, require investigation. Overrides should require supervisory approval.
	\end{frame}
	
	% ============================================================
	% SECTION 14: INTERACTION WITH FRONT OFFICE (RMs)
	% ============================================================
	
	\begin{frame}{Section 6.14: Front Office Interaction - RMs and Compliance}
		\begin{block}{Core Concept}
			Relationship managers (First Line) and compliance (Second Line) have complementary but distinct roles.
		\end{block}
		\textbf{RM Responsibilities:}
		\begin{itemize}
			\item Know your customer (KYC) at onboarding
			\item Identify and report unusual activity to compliance
			\item Respond to RFIs from compliance
		\end{itemize}
		\textbf{Compliance Responsibilities:}
		\begin{itemize}
			\item Investigate suspicious activity independently
			\item Make SAR filing decisions (RMs do NOT decide)
			\item Escalate material findings to management/Board
		\end{itemize}
		\redflag{Critical:} RMs cannot veto or prevent SAR filings based on "commercial importance" of the client.
	\end{frame}
	
	\begin{frame}{6.14.1: The RM Who Blocks a SAR}
		\examfocus{Common Question Type: Commercial Pressure vs. Compliance Obligation}
		
		\scenario{Typical Exam Scenario:}
		A compliance analyst concludes there is reasonable suspicion of money laundering on an account belonging to a high-revenue corporate client. The relationship manager argues: "This client generates \$5 million in annual fees. We can't file a SAR based on suspicion. Get more evidence first." The CCO, under pressure from the RM, delays filing.
		
		\textbf{What the Exam Tests:} RMs cannot override compliance decisions on SAR filings.
		
		\textbf{How to Analyze:}
		\begin{enumerate}
			\item Reasonable suspicion is the legal standard
			\item Commercial importance is NOT a valid reason to delay or avoid filing
			\item The CCO has an independent obligation to file
			\item Delaying filing is a regulatory violation
		\end{enumerate}
		
		\trap{Common Distractor:} "The RM's commercial concern is legitimate and should be considered."
		\correct{Correct Answer Approach:} Compliance decisions on SAR filing must be independent of commercial considerations. The CCO must file upon reasonable suspicion regardless of client revenue.
	\end{frame}
	
	% ============================================================
	% SECTION 15: THIRD-PARTY AND VENDOR DUE DILIGENCE
	% ============================================================
	
	\begin{frame}{Section 6.15: Third-Party and Vendor Due Diligence}
		\begin{block}{Core Concept}
			Banks remain responsible for outsourced AML functions. Vendor risk must be managed.
		\end{block}
		\textbf{Key Requirements:}
		\begin{itemize}
			\item Due diligence on vendors before engagement
			\item Contractual obligations for AML compliance
			\item Right to audit vendor performance
			\item Ongoing monitoring of vendor effectiveness
			\item Bank retains ultimate regulatory responsibility
		\end{itemize}
		\redflag{Cannot outsource responsibility:} Even if a vendor performs AML functions, the bank is still accountable to regulators.
	\end{frame}
	
	\begin{frame}{6.15.1: The Outsourced Monitoring Failure}
		\examfocus{Common Question Type: Vendor Failure and Bank Liability}
		
		\scenario{Typical Exam Scenario:}
		A bank outsources its transaction monitoring to a third-party vendor. The vendor's system fails to generate alerts for 8 months due to a technical error. The bank's internal compliance team did not perform any validation of vendor outputs. When discovered, the vendor claims liability. The bank argues the vendor is responsible.
		
		\textbf{What the Exam Tests:} Banks cannot outsource regulatory responsibility.
		
		\textbf{How to Analyze:}
		\begin{enumerate}
			\item The bank is ultimately responsible to regulators
			\item The bank cannot delegate liability to the vendor
			\item The bank should have performed validation and reconciliation
			\item The bank must remediate (review missed period, file SARs) and address vendor oversight
		\end{enumerate}
		
		\trap{Common Distractor:} "The vendor's contract makes them liable for monitoring failures."
		\correct{Correct Answer Approach:} The bank retains regulatory responsibility regardless of vendor contracts. The bank must have vendor oversight and validation controls.
	\end{frame}
	
	% ============================================================
	% SECTION 16: ESCALATION AND OFFBOARDING
	% ============================================================
	
	\begin{frame}{Section 6.16: Escalation and Offboarding}
		\begin{block}{Core Concept}
			When a customer poses unacceptable risk, the bank must escalate through governance and potentially offboard.
		\end{block}
		\textbf{Escalation Path:}
		\begin{enumerate}
			\item Analyst identifies suspicious activity
			\item Investigation conducted
			\item If SAR filed, notify AML committee
			\item If multiple SARs or high risk, escalate to senior management
			\item Board notification for material risks
			\item Decision to offboard requires senior management approval
		\end{enumerate}
		\redflag{Offboarding must be documented and risk-based, not retaliatory.}
	\end{frame}
	
	\begin{frame}{6.16.1: The Offboarding Decision After Multiple SARs}
		\examfocus{Common Question Type: When to Offboard}
		
		\scenario{Typical Exam Scenario:}
		A bank has filed 7 SARs on an account over 14 months. The account shows clear patterns of structured deposits followed by wires to offshore shell companies. The compliance committee recommends offboarding. The relationship manager argues: "No conviction has occurred. The client is innocent until proven guilty. We should keep the account."
		
		\textbf{What the Exam Tests:} Offboarding does not require a conviction.
		
		\textbf{How to Analyze:}
		\begin{enumerate}
			\item Banks are not courts of law
			\item The standard for offboarding is risk, not conviction
			\item 7 SARs on one account is clear evidence of elevated risk
			\item The bank can (and often should) offboard
		\end{enumerate}
		
		\trap{Common Distractor:} "The customer is innocent until proven guilty, so the bank should keep the account."
		\correct{Correct Answer Approach:} Offboarding is a risk management decision, not a criminal conviction. Multiple SARs justify offboarding regardless of conviction status.
	\end{frame}
	
	% ============================================================
	% CONCLUSION SLIDES
	% ============================================================
	
	\begin{frame}{Domain 6 Summary: Top 20 Most Tested Concepts}
		\begin{enumerate}
			\item \textbf{Three Lines of Defense:} First=Business, Second=Compliance, Third=Audit
			\item \textbf{Independence Violation:} Audit cannot design controls it later tests
			\item \textbf{Board Oversight:} Approves program, holds management accountable, does NOT file SARs
			\item \textbf{False Negatives:} MISSING suspicious activity (worse than false positives)
			\item \textbf{Over-Filtering:} Raising thresholds too high increases false negatives
			\item \textbf{Model Drift:} Models degrade over time; require periodic recalibration
			\item \textbf{Silent Data Failure:} Missing data with no error messages = monitoring gap
			\item \textbf{Alert Adjudication:} Must be documented and consistent across analysts
			\item \textbf{SAR Filing Standard:} Reasonable suspicion (not proof)
			\item \textbf{Tipping-Off:} Disclosing investigation or SAR filing; 314(b) is exception
			\item \textbf{Subpoena Response:} Comply promptly, maintain records, do NOT notify customer
			\item \textbf{"Keep Open" Requests:} Document in writing, get senior management approval
			\item \textbf{De-Risking:} Blanket termination without individual assessment (problematic)
			\item \textbf{Justified Exit:} Individual, documented, risk-based termination (permissible)
			\item \textbf{RFIs:} Must be neutral; cannot disclose investigation
			\item \textbf{Training:} Must be role-based, not one-size-fits-all
			\item \textbf{Whistleblower Protection:} BSA prohibits retaliation for reporting violations
			\item \textbf{Vendor Oversight:} Bank retains responsibility for outsourced functions
			\item \textbf{Offboarding:} Risk decision, does not require conviction
			\item \textbf{Insider Threats:} Monitor employees, segregate duties, enforce mandatory vacation
		\end{enumerate}
	\end{frame}
	
	\begin{frame}{Critical Red Flags - Domain 6 Quick Reference}
		\begin{itemize}
			\item \redflag{Audit designing controls} → Independence violation
			\item \redflag{RM closing alerts without investigation} → First Line overstepping
			\item \redflag{No documentation in case files} → Unverifiable investigation
			\item \redflag{77\% false negative rate on known SARs} → Material deficiency
			\item \redflag{34\% missing transaction data for 14 months} → Data quality failure
			\item \redflag{Teller telling customer "compliance is reviewing"} → Tipping-off
			\item \redflag{Blanket termination of all African banks} → De-risking
			\item \redflag{RM blocking SAR due to "commercial importance"} → Improper pressure
			\item \redflag{Same training for Board and tellers} → Ineffective training
			\item \redflag{PIP issued 3 weeks after internal report} → Potential retaliation
			\item \redflag{Employee override without supervisor approval} → Insider risk
			\item \redflag{Offboarding delayed for "more evidence"} → Misunderstanding of SAR standard
		\end{itemize}
	\end{frame}
	
	\begin{frame}{Exam Day Reminders for Domain 6}
		\begin{itemize}
			\item \textbf{Domain 6 is 30\% of the exam} - second largest domain
			\item \textbf{Three Lines of Defense} is foundational - know who does what
			\item \textbf{Independence} is critical - audit cannot audit its own work
			\item \textbf{SAR standard} = reasonable suspicion, not proof
			\item \textbf{Tipping-off} includes disclosing an investigation exists
			\item \textbf{False negatives} are worse than false positives
			\item \textbf{De-risking} = blanket termination (bad); justified exit = individual assessment (good)
			\item \textbf{Training} must be role-based
			\item \textbf{Whistleblower protections} exist for internal and external reports
			\item \textbf{Banks cannot outsource regulatory responsibility} to vendors
			\item \textbf{Offboarding} is a risk decision, not a conviction
			\item \textbf{RMs cannot override compliance SAR decisions}
		\end{itemize}
		\vspace{0.5cm}
		\centering
		\greennote{\Large Good luck on your CAMS exam!}
	\end{frame}
	
\end{document}